\section*{Contoh}
\addcontentsline{toc}{section}{Contoh}

\begin{contoh}[Atribut synthesized]
Untuk ekspresi \code{2 + 3.5}, tipe anak kiri adalah \code{int} dan
anak kanan \code{float}. Aturan \code{arith(int, float)}
menghasilkan \code{float}. Atribut tipe simpul akar dihitung dari
tipe anak, sehingga merupakan atribut synthesized.
\end{contoh}

\begin{contoh}[Atribut inherited]
Pada blok
\begin{lstlisting}[language=C]
{
    int x = 1;
    {
        int y = x + 1;
    }
}
\end{lstlisting}
lingkungan scope dalam mewarisi tabel simbol scope luar. Atribut
\code{env} pada blok dalam dihitung dari blok luar. Ini adalah
atribut inherited karena nilainya berasal dari orang tua.
\end{contoh}

\begin{contoh}[SDD ke translation scheme]
Aturan SDD
\[
E \to E_1 + E_2 \quad \{ E.type = \text{arith}(E_1.type,
E_2.type) \}
\]
dapat ditulis sebagai aksi di akhir produksi:
\begin{lstlisting}[language=Bison]
expr : expr T_PLUS expr
     { $$.type = type_arith($1.type, $3.type); }
     ;
\end{lstlisting}
Aksi dijalankan saat reduksi, setelah tipe kedua operand diketahui.
\end{contoh}

\begin{contoh}[Dari SDD ke postfix]
Untuk aturan
\[
E \to E_1 + T \quad \{\, E.code = E_1.code \,\|\, T.code \,\|\,
\text{``+''} \,\}
\]
ekspresi \code{(2 + 3) * 4} menghasilkan postfix
\code{2 3 + 4 *}. Atribut \code{code} merupakan atribut synthesized:
nilainya dihitung dari kode anak-anaknya.
\end{contoh}

\begin{contoh}[Dari SDD ke three-address code]
Untuk pernyataan \code{a = b + c * d}, aturan semantik dapat
membangkitkan instruksi
\[
t_1 = c * d, \qquad t_2 = b + t_1, \qquad a = t_2.
\]
Setiap hasil antara disimpan pada nama sementara. Urutan instruksi
mengikuti dependensi pohon: perkalian diselesaikan sebelum
penjumlahan, dan penjumlahan sebelum assignment.
\end{contoh}

\begin{contoh}[Aturan pernyataan pada pohon sintaks]
Untuk program
\begin{lstlisting}[language=C]
int a = 2;
float b = 3.5;
a = b + 1;
\end{lstlisting}
pohon sintaks memuat deklarasi \(a\), deklarasi \(b\), ekspresi
\(b + 1\), dan assignment \(a = b + 1\). Aturan
\(\text{declare}(a,\code{int})\) dan
\(\text{declare}(b,\code{float})\) mengisi tabel simbol. Aturan
\(\text{arith}(\code{float},\code{int})\) menghasilkan
\code{float}. Aturan
\(\text{check\_assign}(\code{int},\code{float})\) menghasilkan
\code{error} karena assignment memerlukan cast.
\end{contoh}
